\section{The Propers: Detailed Commentary}

\subsection{Introit: rejoicing with the angels}
The entrance antiphon calls everyone to rejoice in the Lord. Its celebration
honors all the saints, while its last words name the Son of God as the object
of the angels' praise. Human veneration and angelic praise therefore meet in
Christ. Schuster explains the feast's breadth by naming apostles and martyrs,
clergy and laity, laborers and slaves: the Holy Spirit sanctifies members of
the whole body of Christ. His explanation concerns holiness in every condition
of life, rather than distinctions of prestige.\footnote{Alfredo Ildefonso
Schuster, \work{The Sacramentary}, V (English edition, 1930), pp.~209--210.}

The following Psalm verse, Ps.~32:1, gives the human singers a qualification:
praise befits the upright. Augustine's first exposition identifies uprightness
with submission to God. The contrast is with the person whose will is bent
away from him. Rejoicing in the Lord therefore involves accepting the one
being praised, rather than measuring him by the singer's preference.
Augustine also expressly directs joy toward God rather than oneself.
This gives the feast's opening an account of the act of praise: the just
rejoice because they belong to the Lord whose goodness they acknowledge.
The Psalm citation belongs to this verse; the preceding antiphon is a composed
liturgical text.\footnote{Augustine, \work{Enarrationes in Psalmos}, Ps.~32,
enarratio I, \S1; historical English in NPNF I.8, Psalm XXXIII, \S1.}

\subsection{The Collect and the Sunday commemoration}
The feast's Collect addresses almighty and eternal God as the giver of the
celebration itself. He has granted that the merits of all his saints should be
venerated on one day. The request then concerns the living: the mercy they
desire, sought amid a multitude of intercessors. Merits, intercession, and
mercy are related without being identified. The saints' holiness is honored;
their prayers assist the petitioners; the gift requested comes from God.
Schuster's exposition retains both the one festival and the many pleading for
the faithful.\footnote{Schuster, \work{Sacramentary}, V, p.~210; MR 1962,
p.~702, Collect.}

The second Collect names the congregation as God's people and asks release
from the bonds of sin. The cause confessed is human frailty; the means sought
is God's kindness. Schuster emphasizes the prayer's universal humility. No
member of the people can claim to have no need of pardon. Read beside the
feast's Collect, this prayer distinguishes the veneration of holiness from
the present condition of those offering it. The assembly can celebrate the
saints truthfully while asking God to free it from sins it has itself
contracted.\footnote{MR 1962, p.~416, Collect; Schuster,
\work{Sacramentary}, III (1927), pp.~186--187, commentary on the Collect.}

\subsection{The Lesson: sealing, enumeration, and worship}
Apocalypse 7 interrupts the sequence of threatened judgment. Chapter 6 ends
with the question of who can stand before divine wrath. In chapter 7 an angel
orders the delay of injury to land, sea, and trees until God's servants have
been sealed on their foreheads. The action identifies them as belonging to
God before the vision describes their praise. The appointed Lesson begins
at 7:2; the preceding restraint of the winds is part of its literary context.

John first hears the number of the sealed: 144,000. The repeated twelve
thousand from each tribe occupies a substantial part of the proclamation.
Bede understands the definite total as the completeness of the Church,
descended from the patriarchs either naturally or by imitation of faith.
He appeals to Galatians 3:29: belonging to Christ makes one Abraham's seed.
His arithmetic develops that interpretation: twelve multiplied by twelve,
then expressed in thousands, signifies completeness; the apostolic twelve
likewise preach throughout the world. This is Bede's explanation of the
vision's number, not a prediction of a maximum count of the saved.
\footnote{Bede, \work{The Explanation of the Apocalypse}, trans. Edward
Marshall (1878), I, on Apoc.~7:4, p.~46; Gal.~3:29.}

The next sight exceeds human counting and names nations, tribes, peoples, and
languages. Bede explains the passage as a return to the victory of those who
overcome persecution, following a recapitulation. He also allows the sequence
to distinguish Israel, to whom the Gospel was first preached, from the
salvation of the Gentiles. That alternative keeps a distinction within the
vision which his interpretation of the whole numbered Church need not make
in the same way. Neither explanation erases the text's repeated enumeration.
\footnote{Bede, \work{Explanation}, I, on Apoc.~7:9, p.~52.}

The multitude stands before the throne and the Lamb, dressed in white and
holding palms. Bede identifies the robes with baptism and the palms with the
victory of the Cross, through which Christ's people overcome the world.
He also gives a further explanation of the robes in terms of glory bestowed
by the Spirit. The act of praise confesses salvation to God and the Lamb;
on verse 12 Bede says that the Church acknowledges having received its gifts
from God. The praise of the angels permits two explanations in his commentary:
they can represent the worshipping people, or they can be the heavenly
spirits rejoicing at the salvation of the Gentiles. The latter provides a
particularly direct comparison with the Introit's angels.\footnote{Bede,
\work{Explanation}, I, on Apoc.~7:9--12, pp.~52--53.}

The Lesson ends with the final Amen of verse 12. The elder's explanation of
the white robes and the great tribulation follows in verses 13--17.
Those verses help a reader understand the scene, but their explanatory
dialogue is not proclaimed within this Lesson. The distinction matters
especially when the feast is described as showing the fruits of Christ's
Passion: the appointed palms and the Lamb already permit that discussion,
while the explicit washing of robes in the Lamb's blood belongs to the
surrounding chapter.\footnote{Apoc.~6:12--7:17; Bede,
\work{Explanation}, I, on Apoc.~7:13--17, pp.~53--54.}

\subsection{Gradual: the good promised to those who seek God}
Psalm 33 addresses the saints themselves with a command to fear the Lord.
The title is compatible with continued reverence and dependence. The chant
then promises that those who fear him lack nothing and that seekers of the
Lord will not be deprived of good. It selects verse 10 and the second part
of verse 11. The verse's first statement about the rich lacking and hungering
belongs to the biblical context, not the chant as sung here.

Augustine directly confronts the apparent contradiction. A wicked rich man
can grow old and die wealthy, with an expensive funeral, while a worshipper
who seeks God remains poor. If good means possessions, the promise appears
false. Augustine names justice, truth, charity, faith, and endurance as
riches of the heart. He then identifies Christ as the living bread and
expressly cites the Beatitude of hunger and thirst for justice. His
explanation of the Psalm consequently bears on this Mass's Gospel through a
connection he actually makes.\footnote{Augustine, Ps.~33, enarratio II,
\S\S14--15; NPNF I.8, Psalm XXXIV, \S\S13--14.}

Aquinas preserves several explanations of the same difficulty. Spiritual
treasure, contentment with what suffices, God's providential help, and a
physician's beneficial withholding of food each explain how lack can coexist
with divine care. The physician comparison concerns an apparent deprivation
which serves the patient. It cannot mean that every absence is already
experienced as relief, or that prayer assures immediate material abundance.
The several explanations concern the good God gives, and the way its gift
may differ from what a person initially seeks.\footnote{Thomas Aquinas,
\work{Super Psalmos}, Ps.~33, nn.~10--11, Corpus Thomisticum,
\href{https://www.corpusthomisticum.org/cps31.html}{Psalm 33 commentary};
Latin consulted, paraphrased here.}

\subsection{Alleluia: the invitation to the burdened}
The Alleluia appoints Matthew 11:28. In the chapter, the invitation follows
Christ's thanksgiving for revelation to little ones and his words about the
Father and the Son. It is followed by his invitation to take his yoke and
learn his meekness and humility. Chrysostom includes people burdened by
anxiety, sorrow, and sins. He reads the promise of rest through that
following account of discipleship.\footnote{Matt.~11:25--30;
John Chrysostom, \work{Homilies on Matthew}, XXXVIII.3.}

Chrysostom also asks how an easy yoke agrees with a narrow way and a command
to take up the cross. His answer has several parts: humility already frees
the soul from pride's anxiety; virtue releases it from the heavier bondage of
sin; Christ gives his own example; future glory exceeds present suffering.
He expressly recalls that Christ began the Beatitudes with humility.
His argument thus explains a specific relationship between the Alleluia
and Gospel without making either text promise freedom from suffering.
\footnote{Chrysostom, \work{Matthew}, XXXVIII.3--5.}

\subsection{Gospel: the kingdom promised in the Beatitudes}
Jesus sits on the mountain, his disciples draw near, and he teaches. The
Vulgate order proclaimed here places the meek before those who mourn.
The repeated blessings name dispositions and actions, alongside present
losses and future gifts. Poverty of spirit begins the sequence; the kingdom
is promised again to those persecuted for justice. The final address blesses
disciples reviled falsely for Christ and promises a great heavenly reward.
The Missal ends before the subsequent comparison with the persecuted prophets.

Augustine understands the first blessing as humility and fear of God. The
following Beatitudes describe a connected education: receptiveness to God's
teaching, grief over bondage to lesser things, effort to desire justice,
mercy toward those needing help, purification of the heart, and ordered peace.
The kingdom at the beginning and at the persecution blessing encloses this
formation. He sees the eighth blessing as showing the completeness of what
has been formed, tested by opposition from outside.\footnote{Augustine,
\work{Our Lord's Sermon on the Mount}, I.1.3; I.3.10--I.4.12.}

Chrysostom stresses the freedom with which humility is chosen: social
lowliness imposed by circumstance does not by itself describe the will.
He also insists that mercy concerns actions as well as money, purity cannot
be replaced by generous giving, and peacemaking requires reconciling other
people. The persecution that follows qualifies the praise of peace: agreement
in evil is not the peace Christ blesses. The final blessing similarly
requires both false accusation and fidelity to Christ.\footnote{Chrysostom,
\work{Matthew}, XV.2--3, 6--7.}

\subsection{Offertory: the just in God's hand}
Wisdom 2 presents people who persecute the just man and test his claim to
be God's child. Wisdom 3 answers their judgment of his death. The Offertory
selects and adapts clauses from 3:1--3. Its wording concerns the torment
of malice; the ordinary biblical wording at this point concerns death.
The intervening biblical clauses are not all retained. What the chant
sings is the security and peace of the just whom the unwise suppose
to have perished.\footnote{Wis.~2:12--3:9; MR 1962, p.~703, Offertory.}

In his exposition of Psalm 87, Augustine directly invokes Wisdom 3:1 while
distinguishing Christ's resurrection from the resurrection still awaited by
his members. The righteous dead remain in God's hand while their bodies
await restoration. Christ freely lays down his life and takes it up; his
members receive at the end what preceded them in their Head.
Schuster's explanation of the chant likewise denies the persecutor's
supposed possession of the martyr: the martyr belongs to God.
\footnote{Augustine, Ps.~87, \S5; NPNF I.8, Psalm LXXXVIII, \S4;
Schuster, \work{Sacramentary}, V, p.~211.}

\subsection{The Secrets and the Common Preface}
The feast's Secret concerns both honor paid to all the just and salvation
needed by those making the offering. God's mercy is expressly requested for
that saving benefit. The commemorated Sunday's Secret adds the completion
of an unmerited gift and growth in service. Schuster interprets this service
as the ministry of the altar and the sacrifice of praise through Christ.
The language of offering therefore remains the language of dependence:
the offerers ask God to complete what he has granted.\footnote{MR 1962,
pp.~703, 417, Secrets; Schuster, \work{Sacramentary}, III, p.~188;
V, p.~211.}

The Common Preface gives thanks to the Father through Christ. It names
angels, dominations, powers, heavenly virtues, and seraphim, then asks that
human voices be admitted with them. The assembly's praise has the same
Christological direction as the entrance chant. Its request for admission
also distinguishes the earthly singers' petition from John's description
of the worshippers before the throne. The Preface appointed in this dated
Mass is the Common Preface; the historical Preface printed in Schuster's
feast commentary is a different text.\footnote{MR 1962, pp.~252--253,
Common Preface; RGMR 482--483, 494(b), 498.}

\subsection{Communion and the two Postcommunions}
The Communion repeats three blessings from the Gospel: purity of heart,
peacemaking, and persecution for justice. The selection places the promise
of seeing God beside divine sonship and the kingdom. For Augustine, the
purified heart is single, able to see by the faculty suited to God; peace
requires the person's desires to be ordered under reason and reason under
the Son. Chrysostom adds the practical task of reconciling others.
The communicants hear a description of sanctity that concerns both inward
conversion and the treatment of other people.\footnote{Matt.~5:8--10;
Augustine, \work{Sermon on the Mount}, I.2.8--9;
Chrysostom, \work{Matthew}, XV.6--7.}

The feast's Postcommunion asks for enduring joy in veneration of all the
saints and protection through their continual supplication. The Sunday's
prayer asks that those granted divine participation should not succumb to
human dangers. Schuster explains the latter as a request for strengthening
within the battle, rather than removal from every trial. Participation
and danger coexist in the same petition. The feast thus ends with the living
still seeking help, assisted by the saints they have honored.\footnote{
MR 1962, pp.~703, 417, Postcommunions; Schuster, \work{Sacramentary},
III, p.~189; V, p.~211.}
